% from ega3-1-fr.tex
\oldpage[III]{132}
Sous les hypothèses de (\hyperref[III.4.3.2-fr]{4.3.2}), on peut alors en
renforcer la conclusion~: les fibres $f^{-1}(y)$ sont en effet
\emph{géométriquement connexes}. Pour le voir, observons que pour toute
extension $k$ de $\boldsymbol{k}(y)$, il existe un anneau local noethérien $A$ et un
homomorphisme local $\varphi:\mathcal O_y\to A$ qui fait de $A$ un
$\mathcal O_y$-module \emph{plat} et tel que le corps résiduel de $A$ soit
$\boldsymbol{k}(y)$-isomorphe à $k$ (\textbf{0}, 10.3.1). Soit alors
$Y_1=\operatorname{Spec}(A)$ et soit $h:Y_1\to Y$ le morphisme local
correspondant à $\varphi$, transformant l'unique point fermé $y_1$ de $Y_1$
en $y$ (\textbf{I}, 2.4.1)~; posons
$X_1=X\times_YY_1$, et $f_1=f\times1_{Y_1}$~; $f_1$ est propre
(\textbf{II}, 5.4.2, (iii)) et $f_1^{-1}(y_1)$ est un $\boldsymbol{k}(y_1)$-préschéma isomorphe à
$X\times_Y\operatorname{Spec}(k)$. Tout revient donc à montrer que
$f_{1*}(\mathcal O_{X_1})=\mathcal O_{Y_1}$ pour pouvoir appliquer
(\hyperref[III.4.3.2-fr]{4.3.2}) à $f_1$. Or $g$ est un morphisme \emph{plat},
comme il résulte de (\textbf{I}, 2.4.2) et de
(\hyperref[III.1.4.15-fr]{1.4.15.5})~; on a donc
$f_{1*}(\mathcal O_{X_1})
=h^*(f_*(\mathcal O_X))
=h^*(\mathcal O_Y)
=\mathcal O_{Y_1}$
en vertu de (\hyperref[III.1.4.15-fr]{1.4.15}) appliqué pour $q=0$.

Dans le cas général (\hyperref[III.4.3.1-fr]{4.3.1}), le même raisonnement
montre que l'on a (avec les notations de
(\hyperref[III.4.3.1-fr]{4.3.1}))
$f_{1*}(\mathcal O_{X_1})=h^*(g_*(\mathcal O_{Y'}))$,
et la factorisation de Stein $f_1=g_1\circ f'_1$ de $f_1$ est telle que
$g_1=g\times1_{Y_1}$ (\textbf{II}, 1.5.2), le $Y_1$-schéma fini correspondant étant
$Y'_1=Y'\times_YY_1$. Tenant compte de la transitivité des fibres
(\textbf{I}, 3.6.4), on voit donc que le nombre des composantes connexes de
$f_1^{-1}(y_1)$ est, en vertu de
(\hyperref[III.4.3.3-fr]{4.3.3}), égal au nombre d'éléments de
$g_1^{-1}(y_1)=g^{-1}(y)\otimes_{\boldsymbol{k}(y)}k$.
Si on prend pour $k$ une extension \emph{algébriquement close} de $\boldsymbol{k}(y)$, ce nombre
est indépendant de l'extension algébriquement close considérée et égal au
\emph{nombre géométrique de points de} $g^{-1}(y)$ (\textbf{I}, 6.4.7), ou encore à
la \emph{somme des rangs séparables}
$[\boldsymbol{k}(y'_i):\boldsymbol{k}(y)]_s$
où $y'_i$ parcourt l'ensemble fini $g^{-1}(y)$. On dit encore que ce nombre
est le \emph{nombre géométrique de composantes connexes} de $f^{-1}(y)$. On
notera que les $\boldsymbol{k}(y'_i)$ ne sont autres que les corps résiduels de l'anneau
semi-local $(f_*(\mathcal O_X))_y$.
\end{remark}

\begin{proposition}[4.3.5]
\label{III.4.3.5-fr}
Soient $X$ et $Y$ deux préschémas localement noethériens intègres et
$f:X\to Y$ un morphisme propre dominant. Pour tout $y\in Y$, le nombre de
composantes connexes de $f^{-1}(y)$ est au plus égal au nombre des idéaux
maximaux de la fermeture intégrale $\mathcal O'_y$ de $\mathcal O_y$ dans le
corps des fonctions rationnelles $\mathrm R(X)$.
\end{proposition}

\begin{proof}
En effet, pour tout ouvert $U$ de $Y$,
$\Gamma(U,f_*(\mathcal O_X))=\Gamma(f^{-1}(U),\mathcal O_X)$
est l'intersection des anneaux locaux $\mathcal O_x$ tels que
$x\in f^{-1}(U)$ (\textbf{I}, 8.2.1.1). On en conclut aussitôt que la fibre
$(f_*(\mathcal O_X))_y$ est un sous-anneau de $\mathrm R(X)$ contenant
$\mathcal O_y$. En outre, comme $f_*(\mathcal O_X)$ est un
$\mathcal O_Y$-Module cohérent, $(f_*(\mathcal O_X))_y$ est un
$\mathcal O_y$-module de type fini, donc contenu dans $\mathcal O'_y$~; on
sait ([13], vol. I, p. 257 et 259) que tout idéal maximal d'un tel anneau
$A$ est l'intersection de $A$ et d'un idéal maximal de $\mathcal O'_y$,
d'où la proposition.
\end{proof}

\begin{definition}[4.3.6]
\label{III.4.3.6-fr}
On dit qu'un anneau local intègre est unibranche si sa clôture
intégrale est un anneau local. On dit qu'un point $y$ d'un préschéma intègre
$Y$ est unibranche si l'anneau local $\mathcal O_y$ est unibranche \textup{(ce qui
est en particulier le cas lorsque $Y$ est \emph{normal} au point $y$)}.
\end{definition}

Soit $A$ un anneau local intègre, et soit $K$ son corps des fractions~; pour
que $A$ soit unibranche, il faut et il suffit que tout sous-anneau $A_1$ de
$K$, contenant $A$ et qui est une $A$-algèbre finie, soit un anneau local.
En effet, soit $A'$ la clôture intégrale de $A$~; il résulte du premier
théorème de Cohen-Seidenberg (Bourbaki, \emph{Alg. comm.}, chap. V, \S~2,
n\textsuperscript{o}~1, th.~1) que tout idéal maximal de $A_1$ est la trace
d'un idéal maximal de $A'$, donc si $A'$ est local, il en est de même de
$A_1$. Réciproquement, $A'$ est limite inductive
